\textbf{а)} $\mathbb{F}_4 \cong \mathbb{Z}_2(\alpha): \alpha^2 + \alpha + 1 = 0$ (расширили с помощью корня неприводимого многочлена)

Рассмотрим мультипликативную группу:
$\alpha, \alpha^2 = \alpha + 1, \alpha^3 = 1$. Это циклическая группа на трех элементах, причем, мы не получили только $0$.

Значит, что $\mathbb{F}^*_4 = \mathbb{F}_4 \backslash \{0\} \cong C_3 = \ <\alpha>$\\

\textbf{б)} $\mathbb{F}_8 \cong \mathbb{Z}_2(\alpha): \alpha^3 + \alpha + 1 = 0$

Количество элементов следует из теоремы позапрошлого билета.

$\alpha, \alpha^2, \alpha^3 = \alpha + 1, \alpha^4 = \alpha^2 + \alpha, \alpha^5 = \alpha^2 + \alpha + 1, \alpha^6 = \alpha^2 + 1, \alpha^7 = 1 \Rightarrow <\alpha> \cong \mathbb{F}^*_8$, и мультипликативная группа циклическая\\

\textbf{в)} $\mathbb{F}_9 \cong \mathbb{Z}_3(\alpha): \alpha^2 + 1 = 0$

Рассмотрим $<\alpha + 2>$: $\alpha + 2, (\alpha + 2)^2 = \alpha, (\alpha + 2)^3 = 2\alpha - 1, (\alpha + 2)^4 = -1, (\alpha + 2)^5 = -\alpha - 2, (\alpha + 2)^6 = -\alpha, (\alpha + 2)^7 = -2\alpha + 1, (\alpha + 2)^8 = 1$. Получили циклическую группу на $8$ элементах (однако, если брать $\alpha$ или $\alpha + 1$, то ничего не получится)






